\documentclass[12pt,a4paper]{article}
\usepackage[margin=1in]{geometry}
\usepackage{listings}
\usepackage{xcolor}
\usepackage{graphicx}
\usepackage{amsmath}
\usepackage{amssymb}
\usepackage{hyperref}
\usepackage{fancyhdr}
\usepackage{array}
\usepackage{booktabs}
\usepackage{tikz}

% ============================================================================
% VHDL Code Listing Setup
% ============================================================================
\lstdefinelanguage{VHDL}{
  morekeywords={library, use, entity, port, in, out, inout, is, end, architecture, signal, 
                process, if, then, elsif, else, case, when, others, rising_edge, 
                std_logic, std_logic_vector, begin, all},
  morecomment=[l]{--},
  morestring=[d]{""},
  basicstyle=\ttfamily\small,
  keywordstyle=\color{blue}\bfseries,
  commentstyle=\color{gray}\itshape,
  stringstyle=\color{red},
  breaklines=true,
  showstringspaces=false,
  tabsize=2,
  numbers=left,
  numberstyle=\tiny\color{gray},
  backgroundcolor=\color{lightgray!15},
  frame=single,
}

\lstset{
  language=VHDL,
  basicstyle=\ttfamily\small,
  keywordstyle=\color{blue}\bfseries,
  commentstyle=\color{darkgray}\itshape,
  stringstyle=\color{darkred},
  breaklines=true,
  showstringspaces=false,
  tabsize=2,
  numbers=left,
  numberstyle=\tiny\color{gray},
  backgroundcolor=\color{lightgray!10},
  frame=single,
  xleftmargin=20pt,
}

% Header and Footer
\pagestyle{fancy}
\fancyhf{}
\rhead{VHDL Conditional Statements}
\lhead{Digital System Design}
\cfoot{\thepage}
\lfoot{Semester 6}

% ============================================================================
% Document Start
% ============================================================================
\title{\textbf{VHDL Conditional Statements}\\
        \normalsize Case, If-Else, and When-Else}
\author{Digital System Design Guide}
\date{\today}

\begin{document}

\maketitle
\tableofcontents
\newpage

% ============================================================================
% Section 1: Introduction
% ============================================================================
\section{Introduction to VHDL Conditional Statements}

VHDL provides three primary mechanisms for implementing conditional logic:
\begin{enumerate}
    \item \textbf{Case Statements} - For multiple mutually exclusive conditions
    \item \textbf{If-Else Statements} - For hierarchical conditional logic
    \item \textbf{When-Else Statements} - For concurrent signal assignments
\end{enumerate}

Each statement type has distinct characteristics, use cases, and performance implications in hardware design.

\subsection{When to Use Each Statement}

\begin{table}[h]
\centering
\begin{tabular}{|l|p{4cm}|p{4cm}|}
\hline
\textbf{Statement} & \textbf{Best Used For} & \textbf{Context} \\
\hline
\textbf{Case} & Multiple exclusive conditions with clear priorities & Sequential (process) \\
\hline
\textbf{If-Else} & Hierarchical/nested conditions & Sequential (process) \\
\hline
\textbf{When-Else} & Concurrent signal assignment & Concurrent \\
\hline
\end{tabular}
\end{table}

% ============================================================================
% Section 2: IF-ELSE Statement
% ============================================================================
\section{IF-ELSE Conditional Statement}

\subsection{Basic Structure}

The \texttt{if-else} statement is used within sequential processes to implement hierarchical conditional logic.

\subsubsection{Syntax}
\begin{lstlisting}
if (condition1) then
    -- statements executed if condition1 is TRUE
    
elsif (condition2) then
    -- statements executed if condition2 is TRUE
    -- (only if condition1 was FALSE)
    
else
    -- statements executed if all conditions are FALSE
end if;
\end{lstlisting}

\subsection{Key Characteristics}
\begin{itemize}
    \item Used within \textbf{sequential processes}
    \item Supports \textbf{hierarchical/nested conditions}
    \item Priority is determined by \textbf{order of conditions}
    \item \textbf{Only one branch} executes (no fallthrough)
    \item Can include \textbf{multiple elsif} clauses
\end{itemize}

\subsection{Example 1: D-Latch (IF-ELSE Statement)}

A D-Latch is a level-triggered memory element. When the enable signal is HIGH, the output follows the input. When enable is LOW, the output holds its previous value.

\subsubsection{Truth Table}
\begin{table}[h]
\centering
\begin{tabular}{|c|c|c|c|}
\hline
\textbf{Enable} & \textbf{Data In} & \textbf{Q} & \textbf{Behavior} \\
\hline
1 & 0 & 0 & Output follows input \\
\hline
1 & 1 & 1 & Output follows input \\
\hline
0 & X & Previous & Output holds \\
\hline
\end{tabular}
\end{table}

\subsubsection{VHDL Implementation}

\begin{lstlisting}
library IEEE;
use IEEE.STD_LOGIC_1164.all;

entity D_Latch is
    port (
        Enable  : in  STD_LOGIC;
        Data_In : in  STD_LOGIC;
        Q       : out STD_LOGIC;
        Q_bar   : out STD_LOGIC
    );
end D_Latch;

architecture Behavioral of D_Latch is
    signal Q_internal : STD_LOGIC := '0';
begin
    
    process(Enable, Data_In)
    begin
        if (Enable = '1') then
            -- When Enable is HIGH, output follows input
            Q_internal <= Data_In;
        -- else: Q_internal holds its previous value
        end if;
    end process;
    
    -- Continuous assignment
    Q     <= Q_internal;
    Q_bar <= not Q_internal;
    
end Behavioral;
\end{lstlisting}

\subsubsection{Explanation}
\begin{itemize}
    \item The \texttt{process} sensitivity list includes \texttt{Enable} and \texttt{Data\_In}
    \item When \texttt{Enable = '1'}, the latch becomes transparent (output follows input)
    \item When \texttt{Enable = '0'}, no assignment occurs, so the value is held
    \item This implements level-triggered behavior (latch)
\end{itemize}

\newpage

\subsection{Example 2: D-Flip-Flop (IF-ELSE Statement in Process)}

A D-Flip-Flop is edge-triggered. It samples the input on a clock edge and holds the value until the next clock edge.

\subsubsection{Truth Table}
\begin{table}[h]
\centering
\begin{tabular}{|c|c|c|c|}
\hline
\textbf{CLK} & \textbf{Reset} & \textbf{Data In} & \textbf{Q Output} \\
\hline
$\uparrow$ & 0 & 0 & 0 \\
\hline
$\uparrow$ & 0 & 1 & 1 \\
\hline
0/1 & 1 & X & 0 (Reset) \\
\hline
\end{tabular}

\begin{flushleft}
$\uparrow$ = Rising edge (clock transition from 0 to 1)
\end{flushleft}
\end{table}

\subsubsection{VHDL Implementation}

\begin{lstlisting}
library IEEE;
use IEEE.STD_LOGIC_1164.all;

entity D_FlipFlop is
    port (
        CLK     : in  STD_LOGIC;
        Reset   : in  STD_LOGIC;
        Data_In : in  STD_LOGIC;
        Q       : out STD_LOGIC;
        Q_bar   : out STD_LOGIC
    );
end D_FlipFlop;

architecture Behavioral of D_FlipFlop is
    signal Q_internal : STD_LOGIC := '0';
begin
    
    process(CLK, Reset)
    begin
        -- Asynchronous reset (highest priority)
        if (Reset = '1') then
            Q_internal <= '0';
            
        -- Synchronous data capture (on rising edge of clock)
        elsif rising_edge(CLK) then
            Q_internal <= Data_In;
        end if;
    end process;
    
    Q     <= Q_internal;
    Q_bar <= not Q_internal;
    
end Behavioral;
\end{lstlisting}

\subsubsection{Key Differences from Latch}
\begin{itemize}
    \item Uses \textbf{\texttt{rising\_edge(CLK)}} for edge-triggered behavior
    \item \textbf{Asynchronous reset} has higher priority (checked first)
    \item Only captures data on the \textbf{rising edge} of the clock
    \item Holds data between clock edges
\end{itemize}

\subsubsection{If-Else Hierarchical Priority}
\begin{enumerate}
    \item \textbf{First condition} (\texttt{Reset = '1'}): Always checked first
    \item \textbf{Second condition} (\texttt{elsif rising\_edge(CLK)}): Only evaluated if reset is not asserted
\end{enumerate}

% ============================================================================
% Section 3: CASE Statement
% ============================================================================
\newpage
\section{CASE Statement}

\subsection{Basic Structure}

The \texttt{case} statement selects one of multiple branches based on the value of a single expression.

\subsubsection{Syntax}
\begin{lstlisting}
case (select_signal) is
    when value1 =>
        -- statements executed if select_signal = value1
    
    when value2 =>
        -- statements executed if select_signal = value2
    
    when value3 | value4 =>
        -- statements executed if select_signal = value3 OR value4
    
    when others =>
        -- default case (catches all unspecified values)
end case;
\end{lstlisting}

\subsection{Key Characteristics}
\begin{itemize}
    \item Used within \textbf{sequential processes}
    \item Evaluates a \textbf{single expression}
    \item All possible values must be covered
    \item \textbf{Mutually exclusive} branches (no fallthrough)
    \item Usually synthesizes to \textbf{MUX (Multiplexer)}} logic
    \item More efficient than nested if-elsif-else for many conditions
\end{itemize}

\subsection{Example 1: 4-to-1 Multiplexer (CASE Statement)}

A 4-to-1 multiplexer selects one of 4 input signals based on a 2-bit select signal.

\subsubsection{Block Diagram}
\begin{lstlisting}
    I0 ----\
    I1 -----|---> Y
    I2 -----|
    I3 ----/
              (selected by S[1:0])
\end{lstlisting}

\subsubsection{Truth Table}
\begin{table}[h]
\centering
\begin{tabular}{|c|c|c|}
\hline
\textbf{S[1:0]} & \textbf{Selection} & \textbf{Output Y} \\
\hline
00 & 0 & I0 \\
\hline
01 & 1 & I1 \\
\hline
10 & 2 & I2 \\
\hline
11 & 3 & I3 \\
\hline
\end{tabular}
\end{table}

\subsubsection{VHDL Implementation using CASE}

\begin{lstlisting}
library IEEE;
use IEEE.STD_LOGIC_1164.all;

entity Mux4to1 is
    port (
        I0 : in  STD_LOGIC;           -- Input 0
        I1 : in  STD_LOGIC;           -- Input 1
        I2 : in  STD_LOGIC;           -- Input 2
        I3 : in  STD_LOGIC;           -- Input 3
        S  : in  STD_LOGIC_VECTOR(1 downto 0);  -- 2-bit select
        Y  : out STD_LOGIC            -- Output
    );
end Mux4to1;

architecture Behavioral of Mux4to1 is
begin
    
    process(I0, I1, I2, I3, S)
    begin
        case S is
            when "00" =>
                Y <= I0;
            
            when "01" =>
                Y <= I1;
            
            when "10" =>
                Y <= I2;
            
            when "11" =>
                Y <= I3;
            
            when others =>
                Y <= '0';  -- Default (should not occur)
        end case;
    end process;
    
end Behavioral;
\end{lstlisting}

\subsubsection{Alternative: Concurrent CASE Statement}

For combinational logic, case can be written as a concurrent statement:

\begin{lstlisting}
architecture Behavioral of Mux4to1 is
begin
    
    process(S, I0, I1, I2, I3)
    begin
        case S is
            when "00" => Y <= I0;
            when "01" => Y <= I1;
            when "10" => Y <= I2;
            when "11" => Y <= I3;
            when others => Y <= '0';
        end case;
    end process;
    
end Behavioral;
\end{lstlisting}

\newpage

\subsection{Example 2: 2-to-4 Decoder (CASE Statement)}

A 2-to-4 decoder converts a 2-bit input into 4 mutually exclusive outputs.

\subsubsection{Block Diagram}
\begin{lstlisting}
    I[1:0] ----> [DECODER] ----> O0
                                O1
                                O2
                                O3
\end{lstlisting}

\subsubsection{Truth Table}
\begin{table}[h]
\centering
\begin{tabular}{|c|c|c|c|c|c|}
\hline
\textbf{I[1]} & \textbf{I[0]} & \textbf{O3} & \textbf{O2} & \textbf{O1} & \textbf{O0} \\
\hline
0 & 0 & 0 & 0 & 0 & 1 \\
\hline
0 & 1 & 0 & 0 & 1 & 0 \\
\hline
1 & 0 & 0 & 1 & 0 & 0 \\
\hline
1 & 1 & 1 & 0 & 0 & 0 \\
\hline
\end{tabular}
\end{table}

\subsubsection{VHDL Implementation using CASE}

\begin{lstlisting}
library IEEE;
use IEEE.STD_LOGIC_1164.all;

entity Decoder2to4 is
    port (
        I       : in  STD_LOGIC_VECTOR(1 downto 0);
        Enable  : in  STD_LOGIC;
        O       : out STD_LOGIC_VECTOR(3 downto 0)
    );
end Decoder2to4;

architecture Behavioral of Decoder2to4 is
begin
    
    process(I, Enable)
    begin
        -- Initialize output to 0
        O <= "0000";
        
        if (Enable = '1') then
            case I is
                when "00" =>
                    O <= "0001";  -- Only O0 is HIGH
                
                when "01" =>
                    O <= "0010";  -- Only O1 is HIGH
                
                when "10" =>
                    O <= "0100";  -- Only O2 is HIGH
                
                when "11" =>
                    O <= "1000";  -- Only O3 is HIGH
                
                when others =>
                    O <= "0000";  -- Default
            end case;
        end if;
    end process;
    
end Behavioral;
\end{lstlisting}

\subsubsection{Key Features}
\begin{itemize}
    \item Uses \textbf{case statement} to map each input to unique output
    \item \textbf{Enable signal} controls decoder operation
    \item \textbf{Mutually exclusive outputs} (only one HIGH at a time)
    \item Efficient for address decoding in memory systems
\end{itemize}

% ============================================================================
% Section 4: WHEN-ELSE Concurrent Assignment
% ============================================================================
\newpage
\section{WHEN-ELSE Concurrent Assignment}

\subsection{Basic Structure}

The \texttt{when-else} statement is a concurrent signal assignment used outside of processes.

\subsubsection{Syntax}
\begin{lstlisting}
output_signal <= value1 when condition1 else
                 value2 when condition2 else
                 value3 when condition3 else
                 default_value;
\end{lstlisting}

\subsection{Key Characteristics}
\begin{itemize}
    \item Used for \textbf{concurrent signal assignment}
    \item Evaluates \textbf{multiple independent conditions}
    \item Priority determined by \textbf{order of evaluation}
    \item Must provide \textbf{default value} at the end
    \item Clean, readable syntax for combinational logic
    \item No process required
\end{itemize}

\subsection{Example 1: 4-to-1 Multiplexer (WHEN-ELSE)}

\begin{lstlisting}
library IEEE;
use IEEE.STD_LOGIC_1164.all;

entity Mux4to1_WhenElse is
    port (
        I0 : in  STD_LOGIC;
        I1 : in  STD_LOGIC;
        I2 : in  STD_LOGIC;
        I3 : in  STD_LOGIC;
        S  : in  STD_LOGIC_VECTOR(1 downto 0);
        Y  : out STD_LOGIC
    );
end Mux4to1_WhenElse;

architecture Behavioral of Mux4to1_WhenElse is
begin
    
    -- Concurrent when-else statement
    -- No process required
    Y <= I0 when (S = "00") else
         I1 when (S = "01") else
         I2 when (S = "10") else
         I3;  -- Default (handles "11" and others)
    
end Behavioral;
\end{lstlisting}

\subsubsection{Advantages over CASE}
\begin{itemize}
    \item \textbf{No process} required (concurrent)
    \item More \textbf{concise} for simple conditions
    \item Assigns directly to output signal
    \item Cleaner for small selector logic
\end{itemize}

\subsection{Example 2: 2-to-4 Decoder (WHEN-ELSE)}

\begin{lstlisting}
library IEEE;
use IEEE.STD_LOGIC_1164.all;

entity Decoder2to4_WhenElse is
    port (
        I       : in  STD_LOGIC_VECTOR(1 downto 0);
        Enable  : in  STD_LOGIC;
        O       : out STD_LOGIC_VECTOR(3 downto 0)
    );
end Decoder2to4_WhenElse;

architecture Behavioral of Decoder2to4_WhenElse is
begin
    
    -- Concurrent when-else statement
    O <= "0001" when (Enable = '1' and I = "00") else
         "0010" when (Enable = '1' and I = "01") else
         "0100" when (Enable = '1' and I = "10") else
         "1000" when (Enable = '1' and I = "11") else
         "0000";  -- Default (disabled or invalid)
    
end Behavioral;
\end{lstlisting}

\subsection{Example 3: Serial-to-Parallel Shift Register Output}

WHEN-ELSE can be used for logic involving the shift register:

\begin{lstlisting}
library IEEE;
use IEEE.STD_LOGIC_1164.all;

entity ShiftRegister_Control is
    port (
        Shift       : in  STD_LOGIC;
        Load        : in  STD_LOGIC;
        Serial_In   : in  STD_LOGIC;
        Reset       : in  STD_LOGIC;
        Mode_Select : in  STD_LOGIC_VECTOR(1 downto 0);
        Output      : out STD_LOGIC
    );
end ShiftRegister_Control;

architecture Behavioral of ShiftRegister_Control is
begin
    
    -- Priority-based control logic using when-else
    Output <= '0' when (Reset = '1') else
              Serial_In when (Load = '1') else
              '1' when (Shift = '1') else
              '0';  -- Default idle state
    
end Behavioral;
\end{lstlisting}

% ============================================================================
% Section 5: Comparison Table
% ============================================================================
\newpage
\section{Comparison: IF-ELSE vs CASE vs WHEN-ELSE}

\begin{table}[h]
\centering
\begin{tabular}{|l|c|c|c|}
\hline
\textbf{Feature} & \textbf{IF-ELSE} & \textbf{CASE} & \textbf{WHEN-ELSE} \\
\hline
Context & Sequential & Sequential & Concurrent \\
\hline
Within Process & Yes & Yes & No \\
\hline
Hierarchy & Supports nesting & Flat evaluation & Flat (priority order) \\
\hline
Best for & Complex conditions & Multiple values & Simple multiplexing \\
\hline
Priority & Top to bottom & Mutual exclusive & Top to bottom \\
\hline
Synthesizes to & Logic gates & Multiplexer & Multiplexer \\
\hline
Default required & No & \texttt{when others} & Yes \\
\hline
Readability & Good & Excellent & Good \\
\hline
Performance & Varies & Optimal & Optimal \\
\hline
\end{tabular}
\end{table}

% ============================================================================
% Section 6: Selection Guide
% ============================================================================
\newpage
\section{Selection Guide}

\subsection{Use IF-ELSE When:}
\begin{itemize}
    \item You have \textbf{hierarchical or nested conditions}
    \item Conditions are \textbf{complex} or interdependent
    \item You need \textbf{sequential logic} (latches, flip-flops)
    \item Some conditions have \textbf{higher priority} than others
    \item Working within a \textbf{process block}
\end{itemize}

\textbf{Example:} D-Latch, D-Flip-Flop, or reset logic

\subsection{Use CASE When:}
\begin{itemize}
    \item Selecting from \textbf{multiple discrete values}
    \item Values are \textbf{mutually exclusive}
    \item Need \textbf{clean, readable code}
    \item Working with \textbf{multi-bit selectors}
    \item Inside a \textbf{process block}
\end{itemize}

\textbf{Example:} Multiplexers, Decoders, State machines

\subsection{Use WHEN-ELSE When:}
\begin{itemize}
    \item Implementing \textbf{combinational logic} only
    \item Want \textbf{concurrent assignment}
    \item Conditions are \textbf{simple and independent}
    \item Don't need a \textbf{process block}
    \item Code is \textbf{short and clear}
\end{itemize}

\textbf{Example:} Simple multiplexers, conditional routing, control signals

% ============================================================================
% Section 7: Advanced Example - Serial-to-Parallel Shift Register
% ============================================================================
\newpage
\section{Advanced Example: 4-Bit Serial-to-Parallel Shift Register}

\subsection{Overview}

A 4-bit Serial-to-Parallel (S2P) shift register demonstrates hierarchical if-else logic combined with case statements for control.

\subsection{Complete VHDL Implementation}

\begin{lstlisting}
library IEEE;
use IEEE.STD_LOGIC_1164.all;

entity Serial_to_Parallel_SR is
    port (
        CLK         : in  STD_LOGIC;
        RESET       : in  STD_LOGIC;
        Shift       : in  STD_LOGIC;
        Serial_In   : in  STD_LOGIC;
        S           : out STD_LOGIC_VECTOR(3 downto 0);
        Control     : in  STD_LOGIC_VECTOR(1 downto 0)
    );
end Serial_to_Parallel_SR;

architecture Behavioral of Serial_to_Parallel_SR is
    signal S_reg : STD_LOGIC_VECTOR(3 downto 0);
    
begin
    
    -- ====================================================
    -- MAIN PROCESS: Using IF-ELSE for hierarchical logic
    -- ====================================================
    process(CLK, RESET)
    begin
        -- HIGHEST PRIORITY: Asynchronous Reset
        if (RESET = '1') then
            S_reg <= (others => '0');
        
        -- SECOND PRIORITY: Clock edge check
        elsif rising_edge(CLK) then
            
            -- CASE statement for shift control
            case Shift is
                when '1' =>
                    -- Shift enable: shift data right
                    S_reg <= Serial_In & S_reg(3 downto 1);
                
                when '0' =>
                    -- Shift disable: hold value
                    S_reg <= S_reg;
                
                when others =>
                    S_reg <= (others => '0');
            end case;
        end if;
    end process;
    
    -- ====================================================
    -- CONCURRENT ASSIGNMENT: Using WHEN-ELSE
    -- for output multiplexing based on Control signal
    -- ====================================================
    S <= S_reg when (Control = "00") else
         not S_reg when (Control = "01") else
         (others => '0') when (Control = "10") else
         (others => '1');  -- Default
    
end Behavioral;
\end{lstlisting}

\subsection{Analysis}

\textbf{Hierarchical IF-ELSE Structure:}
\begin{enumerate}
    \item \textbf{First condition:} \texttt{RESET = '1'} (highest priority)
    \item \textbf{Second condition:} \texttt{rising\_edge(CLK)} (synchronous operation)
    \item \textbf{Inside CLK block:} CASE statement for shift control
\end{enumerate}

\textbf{Why Each Statement Type:}
\begin{itemize}
    \item \textbf{IF-ELSE:} Handles priority (reset > clock), asynchronous to synchronous transitions
    \item \textbf{CASE:} Manages shift control (on/off selection)
    \item \textbf{WHEN-ELSE:} Concurrent output routing based on mode control
\end{itemize}

% ============================================================================
% Section 8: Best Practices
% ============================================================================
\newpage
\section{Best Practices}

\subsection{General Design Guidelines}

\begin{enumerate}
    \item \textbf{Always cover all cases}
        \begin{itemize}
            \item Add \texttt{when others} to case statements
            \item Provide default value in when-else
            \item Don't leave conditions undefined
        \end{itemize}
    
    \item \textbf{Use meaningful signal/variable names}
        \begin{itemize}
            \item \texttt{Enable} instead of \texttt{E}
            \item \texttt{Serial\_In} instead of \texttt{SI}
            \item Make intent clear in code
        \end{itemize}
    
    \item \textbf{Maintain hierarchical structure}
        \begin{itemize}
            \item Keep reset logic separate (highest priority)
            \item Use if-elsif-else for dependent conditions
            \item Nest case inside if blocks when needed
        \end{itemize}
    
    \item \textbf{Avoid implicit logic}
        \begin{itemize}
            \item Don't rely on signal holding values implicitly
            \item Explicitly assign default values
            \item Make sequential logic explicit with rising\_edge()
        \end{itemize}
    
    \item \textbf{Choose the right tool}
        \begin{itemize}
            \item Use case for selectors (clean, synthesizes well)
            \item Use if-else for complex conditions
            \item Use when-else for simple concurrent logic
        \end{itemize}
\end{enumerate}

\subsection{Synthesis Considerations}

\begin{table}[h]
\centering
\begin{tabular}{|l|p{5cm}|p{4cm}|}
\hline
\textbf{Pattern} & \textbf{Synthesis Result} & \textbf{Notes} \\
\hline
Simple case with few branches & AND-OR gates & Optimal \\
\hline
Case with many (4+) branches & Multiplexer trees & Efficient \\
\hline
Nested if-elsif-else & Cascaded logic & May have long delays \\
\hline
When-else (priority) & Multiplexer chain & Fast for small selectors \\
\hline
Asynchronous reset & Set/Reset logic & Built into flip-flops \\
\hline
Multiple elsif branches & Priority detection logic & Priority is implicit \\
\hline
\end{tabular}
\end{table}

% ============================================================================
% Section 9: Common Mistakes to Avoid
% ============================================================================
\newpage
\section{Common Mistakes and How to Avoid Them}

\subsection{Mistake 1: Missing WHEN OTHERS}

\textbf{❌ Wrong:}
\begin{lstlisting}
case Mode is
    when "00" => Output <= A;
    when "01" => Output <= B;
    when "10" => Output <= C;
    -- Missing: when others => ...
end case;
\end{lstlisting}

\textbf{✓ Correct:}
\begin{lstlisting}
case Mode is
    when "00" => Output <= A;
    when "01" => Output <= B;
    when "10" => Output <= C;
    when others => Output <= '0';  -- Always provide default
end case;
\end{lstlisting}

\subsection{Mistake 2: Incomplete Sensitivity List}

\textbf{❌ Wrong:}
\begin{lstlisting}
process(CLK)  -- Missing: Reset
begin
    if (Reset = '1') then  -- May not work!
        Q <= '0';
    elsif rising_edge(CLK) then
        Q <= D;
    end if;
end process;
\end{lstlisting}

\textbf{✓ Correct:}
\begin{lstlisting}
process(CLK, Reset)  -- Include all async signals
begin
    if (Reset = '1') then
        Q <= '0';
    elsif rising_edge(CLK) then
        Q <= D;
    end if;
end process;
\end{lstlisting}

\subsection{Mistake 3: Confusing Priority}

\textbf{❌ Problem:} WHEN-ELSE evaluates top-to-bottom
\begin{lstlisting}
Y <= B when (Enable = '0') else  -- Checked FIRST
     A when (Enable = '1') else   -- Checked SECOND (too late!)
     A;
\end{lstlisting}

\textbf{✓ Correct:}
\begin{lstlisting}
Y <= A when (Enable = '1') else   -- Checked FIRST
     B when (Enable = '0') else
     A;  -- Default
\end{lstlisting}

\subsection{Mistake 4: Invalid Use of CASE Outside Process}

\textbf{❌ Wrong:} (CASE cannot be concurrent)
\begin{lstlisting}
architecture Bad of Example is
begin
    -- This is INVALID - case is not concurrent
    case S is
        when "00" => Output <= A;
        when others => Output <= B;
    end case;
end Bad;
\end{lstlisting}

\textbf{✓ Correct:}
\begin{lstlisting}
architecture Good of Example is
begin
    process(S, A, B)
    begin
        case S is
            when "00" => Output <= A;
            when others => Output <= B;
        end case;
    end process;
end Good;
\end{lstlisting}

% ============================================================================
% Section 10: Summary
% ============================================================================
\newpage
\section{Summary}

\subsection{Key Takeaways}

\textbf{IF-ELSE:}
\begin{itemize}
    \item Hierarchical conditional logic
    \item Priority determined by order
    \item Best for complex conditions and latches/flip-flops
    \item Use inside processes
\end{itemize}

\textbf{CASE:}
\begin{itemize}
    \item Select from multiple values
    \item Mutually exclusive branches
    \item Synthesizes efficiently to multiplexers
    \item Always include \texttt{when others}
\end{itemize}

\textbf{WHEN-ELSE:}
\begin{itemize}
    \item Concurrent signal assignment
    \item Priority by order of conditions
    \item Clean for simple combinational logic
    \item Must have default value
\end{itemize}

\subsection{Quick Reference}

\begin{table}[h]
\centering
\begin{tabular}{|l|c|c|c|}
\hline
\textbf{Need...} & \textbf{Use IF-ELSE} & \textbf{Use CASE} & \textbf{Use WHEN-ELSE} \\
\hline
Sequential logic & ✓ & ✓ & ✗ \\
\hline
Combinational logic & ✓ & ✓ & ✓ \\
\hline
Nested conditions & ✓ & ✗ & ✗ \\
\hline
Multiple selectors & ✗ & ✓ & ✓ \\
\hline
Concurrent code & ✗ & ✗ & ✓ \\
\hline
Asynchronous reset & ✓ & ✓ & ✗ \\
\hline
\end{tabular}
\end{table}

% ============================================================================
% References
% ============================================================================
\newpage
\begin{thebibliography}{99}

\bibitem{ieee} IEEE Std 1076-2019: VHDL Language Reference Manual

\bibitem{vhdl_guide} Perry, Douglas L. (2002). \textit{VHDL Programming by Example}. 
McGraw-Hill, 4th Edition.

\bibitem{synthesis} Mealy, George H. \& Blakeslee, Thomas R. (1997). \textit{Digital Logic 
Design and Computer Organization}. John Wiley \& Sons.

\bibitem{dsd} Mano, M. Morris \& Ciletti, Michael D. (2013). \textit{Digital Design}. 
Pearson, 5th Edition.

\end{thebibliography}

\end{document}
